\section{Canonical laminations}

\subsection{Ideally marked hyperbolic plane}
 

\sss{Quadrilaterals}
  Let $\BPi$ be a  standard rectangle with horizontal sides $\BI$ and
  $\BJ$. It supports  the vertical foliation $\FF$. We let 
$$
   \WW(\BI, \BJ) =  \WW(\BPi)  = \WW(\FF) 
$$
be its conformal {\it width}, which is also equal to the width of the 
{\it full path family} in $\Pi$ connecting $\BI$ to $\BJ$.  

\begin{rem}
  We will sometimes use notation $\WW\{ \BI, \BJ\}$ to emphasize that the
  pair of intervales $\BI$ and $\BJ$ is non-ordered. 
\end{rem}

Let $\Bga$ and $\Bde$ be the hyperbolic geodesics in $\Pi$ sharing
the endpoints with $\BI$ and $\BJ$ respectively. We let $d_\hyp\equiv d_\Pi$ be the
hyperbolic distance in $\BPi$, and $l_\hyp\equiv l_\BPi$ be the
 hyperbolic length  (similar notation will be used for other
hyperbolic Riemann surfaces). 

\begin{lem}\label{geod in quadr}
    Let $\eps>0$ and let $\Bga^\eps=\{ z\in\Bga: \ \dist_\hyp (z, \Bde)
    < \eps\} $. Then 
$$
       l_\hyp (\Bga^\eps) = \WW(\BPi) + O_\eps (1).
$$
\end{lem}

Let us now formulate an important transformation rule: 

\begin{lem}\label{tranform rule for rectangles}
   Let $\Bi: \BPi'\ra \BPi$ be a holomorpic map between two rectangles
which extends to homeomorphisms between the respective  horizontal  sides
$\BI'\ra \BI$ and $\BJ'\ra \BJ$. Then
$$
      \WW(  \BPi' ) \leq \WW ( \BPi ). 
$$ 
\end{lem}

\begin{proof}
Let $\De$ and $\De'$ be the horizontal path families in $\BPi$
and $\BPi'$ respectively.  
 Any  path  $\de$ in $\BPi$  contains a subpath that lifts to a 
  path in $\De'$.  Hence 
$$
     \WW(\BPi)= \LL(\De) \geq \LL(\De') = \WW(\BPi').
$$
\end{proof}

\sss{Canonical lamination}
 
  Let us consider the hyperbolic plane in the disk model $\D$. 
An {\it ideal marking } of $\D$ is a {\em tiling}  $\II$
of the circle  $\T=\di \D$ into intervals $\BI_1, \dots \BI_p$
(written in the order they appear on $\T$, up to a cyclic permutation). 
Thus, we can label the intervals by elements of $\Z/ p\Z$. 
We call $p$ the {\it cardinality } of the marking. 

Any non-adjacent pair of intervals $\BI_m$, $\BI_n$
determines an ideal quadrilateral $\ol \BPi_{mn}$ with ``horizontal'' sides $\BI_m $ and  $\BI_n $.
Let $\ol \FF_{mn}$ be the vertical harmonic foliation in $\ol \BPi_{mn}$,
and let  
$$
   \ol \WW_{mn}\equiv \ol \WW(\BI_m, \BI_n) \equiv \WW (\ol  \BPi_{mn} ) \equiv \WW(\ol \FF_{mn} )
$$ 
be the width of this quadrilateral.
It is equal to the width of the full path family in $\D$  connecting $\BI_n$ to $\BI_m$. 

Removing  the square buffers from the  quadrilaterals $\ol \BPi_{mn}$
with width  greater than 2,
we obtain quadralaterals $\BPi_{mn}$ supporting foliations $\FF_{mn}$
with widths
$$
  \WW_{mn} \equiv  \WW(\BI_m, \BI_n) \equiv \WW (\BPi_{mn}  ) \equiv
  \WW(\FF_{mn} )  = \ol \WW_{mn} -2. 
$$

\begin{rem}
   Again, we sometimes wlll use notation  $\WW_{\{ mn\} } \equiv
   \WW\{\BI_m, \BI_n\}$ to emphasize that the pair $\{m,n\}$ is not
   ordered. 
\end{rem}

\begin{lem}\label{canon fol}
    The foliations $\FF_{mn}$ together  form a lamination
    $\FF(\D,\II)$   on $\D$ (supported on $\bigsqcup \BPi_{mn}$).  
\end{lem}

\begin{proof}
  Follows from the Non-Intersection Principle.
\end{proof}

The lamination $\FF(\D,\II)$ is called the {\it canonical lamination
of the  marked hyperbolic plane}. 

\begin{cor}
   There exits at most $p-2$ quadrilaterals $\BPi_{\{mn\}} $.
\end{cor}

\sss{Conformal geometry vs hyperbolic}

Let $\Bga_n$ be the hyperbolic geodesic in $(\D, \II)$ sharing the
endpoints with the $\BI_n$.
Let
$$
\Bga^\eps_{mn}= \{ z\in\Bga_m : \ \dist_\hyp (z, \Bga_n) < \eps\}.
$$
%
% For non-adjacent $I_m$ and $I_n$ we let
% $$
%      d_{mn}\equiv d(\geoD_m, \geoD_n)
% $$
% be the hyperbolic distance between $\geoD_m$ and $\geoD_n$. 
%
% \sss{Thin part} 

Lemma \ref{geod in quadr}  implies: 

\begin{lem}\label{hyp d vs length}
   We have: $\WW_{mn} = l_\hyp(\Bga_{mn}^\eps )   + O_\eps (1). $
\end{lem}

\sss{Thick part}

let $\BGa= \cup \Bga_n$. 
It is  the polygonal curve that bounds  the convex hull $\BP$  of $\di \II= \bigcup \di \BI_n$.  
Let $ \BGa^\eps = \bigcup \Bga^\eps_{mn}$ be the $\eps$-{\it thin part} of $\BGa$. 

\begin{prop}\label{thick part estimate}
 For an ideally  marked hyperbolic plane $(\D, \II)$ with $p$
 intervals $I_n$ and for any $\eps>0$, we have:
$$
       l_\hyp (\BGa \sm \BGa^\eps) = O(p/\eps).
$$
\end{prop}

\begin{proof}
  Let $L= l_\hyp (\BGa \sm \BGa^\eps) $.
% and let $P$ be the ideal polygon
%  bounded by $\BGa$  (i.e., $P$ is the convex hull of $\di\II$). 
  Let us consider a biggest net of $\eps$-separated points $\{ x_i\}$ in $\BGa\sm
  \BGa^\eps$. It contains approximately  $ L/\eps$ points. The
  hyperbolic  disks $\BD_i$  of
  radius $\eps/2$ centered at these points do not overlap, so the total
  hyperbolic area of the half-disks $\BD_i \cap \BP$  is bouded by the hyperbolic area of the polygon
  $\BP $, which is equal to  $\pi (p-2)$. It gives
$$
        \frac L\eps \cdot \eps^2  \leq C  p.
$$
\end{proof}


 \sss{Thin-thick Decomposition}

Let $\ol  \WW_n \equiv \bar \WW(\BI_n)$ be the width of the path family connecting $\BI_n$ 
all other non-adjacent intervals (i.e., to  $\displaystyle{ \T\sm \bigcup_{|k-n| >1} \BI_k } $),
and let 
$$
           \WW_n\equiv    \WW(\BI_n) =   \max  \{ \ol \WW(\BI_n)  -2, \ 0 \} .
$$ 
We call
$$
      \BW \equiv \BW(\D, \II ) :  = \sum \WW_n 
$$
the {\it total weight} of the marked hyperbolic plane. 

\begin{thm}\label{main estimate} 
For any ideal marking $\II = (\BI_1, \dots, \BI_p)$ of the hyperbolic plane
$\D$,  we have:
$$
         \BW (\D, \II)   - O(p)    \leq   2 \sum \WW\{\BI_n, \BI_m\}  
     \leq    \BW(\D, \II ) 
$$
$($where the summation in the middle is taken over non-ordered pairs
$\{ \BI_m, \BI_n\}$$)$. 
\end{thm} 

\begin{proof}
  Since the rectangles $\BPi_{mn}$ are disjoint, we have for each $m$: 
$$
    \sum_n \WW_{\{mn\}} \leq \WW_m.
$$
Summation over $m$ gives the right-hand side estimate (note that each
$\WW_{ \{mn\} }$ is counted twice as we sum  them  up). 

\msk To prove the left-hand estimate, let us take $\eps$ to be the
Margulis constant. Then the  geodesic intervals $\Bga^\eps_{mn}$ are
pairwise disjoint.  Let 
$$
   \BGa^\thin = \bigsqcup_{|m-n|>1} \Bga_{mn}^\eps 
$$
(where $|m-n|$ means the natural distance in $\Z/ p\Z$ such that $m$
and $m+1$ stay distance 1 apart).  
Together with Lemma \ref{hyp d vs length} this implies:

$$
     2\sum \WW_{ \{mn\} } \geq \sum_{|m-n|>1} l_\hyp (\Bga_{mn}^\eps ) - O(p)
     =  l_\hyp ( \BGa^\thin  ) - O(p).
$$

Combining this with Lemma \ref{thick part estimate}, we obtain: 
\begin{equation}\label{thin-thick part}
     2\sum \WW_{ \{mn\} } \geq   l_\hyp ( \BGa^\thin \cup \BGa^\thick ) - O(p).
\end{equation}


Let us now consider the cuspidal  part of $\BGa$:  
$$
 \BGa^\cusp=  \bigcup_{|m-n|=1} \Bga^\eps_{mn}.
$$
We have a partition: 
$$
   \BGa =    \BGa^\thin \cup \BGa^\cusp \cup \BGa^\thick. 
$$
Incorporating it into (\ref{thin-thick part}), we obtain:
$$
    2\sum \WW_{ \{mn\} } \geq   l_\hyp ( \BGa\sm \BGa^\cusp ) - O(p).
$$

Finally, let $\Bga_m^\eps$ consists of those points $z\in \Bga_m$ that
stay distance less than $\eps$ from the geodesic connecting the
endpoints of $\Bga_{m-1}\cup \Bga_m \cup \Bga_{m+1}$. Since $\eps$ is
the Margulis constant, $\Bga_m^\eps \subset \Bga_m\sm \BGa^\cusp$, 
so the last estimate implies
$$
       2\sum \WW_{ \{mn\} } \geq  \sum  l_\hyp ( \Bga_m^\eps ) - O(p).
$$
By Lemma \ref{hyp d vs length},
$ l_\hyp ( \Bga_m^\eps ) = \WW_m + O(1)$, and we are done.
\end{proof}

\sss{Local weights and fluxes for a partially marked plane}

Let $\II= (\BI_1, \dots, \BI_p)$ be a disjoint family of closed
intervales in $\T$.
% with disjoint interiors.
 It is called a {\it partial marking} of $\D$ (of cardinality $p$). 
The {\it gaps} $G_k$  of a partially marked plane are the closures of the
components of $\displaystyle{ \T \sm \bigcup \BI_n}$. Of course, putting
together the intervals and the gaps, we obtain a full ideal  marking
of $\D$. The {\it local weights} $\WW(\BI_n)$ and {\it fluxes}
$\WW(\BG_k)$ are understood in the sense of this full marking.     
 
Let us consider two partially marked copies of the hyperbolic planes, $(\D, \II)$ 
and $ (\D', \II')$, of the same cardinality. We write 
$\Bi : (\D',\II')  \ra (\D, \II) $ for a map $\Bi: \D' \ra \D$
that  extends to a  homeomorphism   $\BI_n' \ra \BI_n$ for any  $n=1,\dots, p$.   
In this case, we have an obvious consistent labeling  of the gaps
$\BG_k$ and $\BG_k'$. 

\begin{lem}\label{transform rules on D}
  For any holomorphic map $\Bi: (\D', \II') \ra (\D, \II) $ between
  partially marked hyperbolic disks we have:

\ssk\nin $\mathrm { (i) } $
$
   \WW(\BI_m' , \BI_n')  \leq \WW(\BI_m, \BI_n). 
$

\ssk\nin $ \mathrm {(ii)} $  $\WW(\BG_k' ) \geq \WW(\BG_k)$. 
\end{lem}

\begin{proof}
  The first estimate follows directly from 
Lemma \ref{tranform rule for rectangles} .  For the second one, observe that
that $\WW (\BG_k)= \WW (\BI_m, \BI_n )^{-1}$, where $\BI_m$ and
$\BI_n$ are the marked intervals adjacent to the gap $\BG_k$.  
\end{proof}

\subsection{Riemann surfaces with boundary}

\sss{Universal covering}

Let us now consider a compact Riemann surface $S$ with non-empty
boundary $\di S = \bigsqcup J_k$. We let $\pi: \D \ra \inter S$ be the 
universal covering, 
and we let $\La\subset \T$ be the limit set for the  group $\De$
of deck transformations.  Then $\pi$ extends continuously to $\T\sm
\La$, and
$$
    \pi:  \ol \D \sm \La \ra  S 
$$ 
is the  universal covering of $S$. Moreover, $\pi$ restricted to  any component of
$\T\sm \La$  gives us a universal covering of some component $J_k$ of
$\di S$.  Such a component of $\T\sm \La$ will be denoted $\BJ_k$
(usually we will need only one component for each $k$, 
so we will not use an extra label for it). 
The stabilizer of $\BJ_k$ in $\De$ is a cyclic groups generated by
a hyperbolic M\"obius transformation $T_k$ with fixed points in  $\di \BJ_k$.


\sss {Local weghts}

We let $\A_k$ be the covering annulus of $S$ corresponding to the
cyclic group $<T_k>$. We call its width $\WW(\A_k)$ 
the {\it local  weight} of $J_k$. We will also use notation $\WW(J_k)$
for this weight.  

 We let
$$
  \BW(S) = \sum \WW(J_k)
$$
be the {\it total weight} of $S$. 

% \begin{lem}
%    We have: $\WW(J_k) = \WW)\BJ_k + O(1) $.
% \end{lem}

\sss{Canonical arc diagram}
  Here we will summarize some notions and results  from \cite{K}. 

  An {\it arc}  $\alpha$ is a proper path in $S$ up to proper homotopy. 
Any arc connects some components $J_k$ and $J_i$ of $\di S$
and lifts to an arc $\ba$ in $\D$ connecting some intervals $\BJ_k$
and $\BJ_i$. We let  $\ol\WW(\alpha)= \ol\WW(\BJ_k, \BJ_i)$.
It is the width of the quadrilateral $\ol \BPi_\alpha=\ol \D$ with
horizontal sides $\BJ_k$ and $\BJ_i$.  
It is independent of the lift used. 

If $\ol\WW(\alpha)> 2$ then
by removing from $\ol\BPi_\alpha$ the square buffers we obtain a
rectangle $\BPi_\ba$.  In this case we let
$$
      \WW(\alpha) = \WW(\BPi_\ba)  = \ol \WW(\alpha)-2.
$$ 
If $\bar\WW(\alpha)\leq 2$ we let  $\WW(\alpha) =0$.

Non-degenerate quadrilaterals $\BPi_\alpha$ project injectively to 
disjoint  {\it canonical quadrilaterals} $\Pi_\alpha\subset S$.  
Each of them supports the canonical
vertical foliation. Altogether these foliations form the
{\it canonical lamination} of $S$. 

Arcs with positive weight form the {\it canonical arc diagram}
$\AAA_S$ on $S$.

% \sss{Thick surfaces}

\sss{Thin-thick decomposition}

\begin{thm}\label{thin-thick for S}
   Let $S$ be  a compact Riemann surface with boundary,
$\chi=\chi(S)$ be its Euler characteristic. Then  we have:
$$
    \BW(S)  -  O( | \chi | ) \leq   2 \sum_{\alpha\in \AAA_S}  \WW(\alpha)  \leq \BW(S). 
$$
\end{thm}

\begin{proof}
    For any component $J_k$ of $\di  S $, $\WW(J_k)\geq \sum
    \WW(\alpha)$, where the summation is taken over the arcs that land
    on $J_k$. This implies that right-hand side estimate. 

 To prove the other one, let us consider the simple closed geodesics
 $\gamma_k$ homotopic to the $J_k$. Then 
\begin{equation}\label{geod and annuli} 
        \WW_k =  \frac 1\pi\,   l_\hyp(\gamma_k) \quad \mathrm{so}\ 
\BW  (S) = \frac 1\pi \sum l_\hyp(\gamma_k). 
\end{equation}

Let $P$ be the Riemann surface bounded the geodesics $\gamma_k$. 
As in the proof of Proposition \ref{thick part estimate}, let $\eps$ be the Margulis
constant for $P$ (equal to the Margulis constant for the double of $P$).

Take some arc $\alpha$ connecting $J_k$ to $J_i$, and lift it to the
universal covering. We obtain an arc $\ba$ in $\D$ connecting $\BJ_k$
to $\BJ_i$. Moreover, the geodesics $\gamma_k$ and $\gamma_i$ lift to
the geodesics $\Bga_k$ and $\Bga_i$ in $\D$ sharing the endpoints with $\BJ_k$ and $\BJ_i$
respectively. By Lemma \ref{geod in quadr}, 
\begin{equation}\label{rectangle  and geod}   
   \WW(\alpha) \equiv \WW(\ba) =\frac 1\pi\, l_\hyp (\Bga_{ki} ^\eps)
   + O(1)= \frac 1\pi \, l_\hyp (\gamma_{ki} ^\eps) +O(1),
\end{equation}
where $\gamma_{ki}^\eps$ are defined on $S$ in the same way as the
$\Bga_{ki}^\eps$ are defined in $\D$, and in fact, the latter is the
lift of the former.  

Let 
$$
   \gamma_k^\thick: =  \gamma_k \sm \bigsqcup_i \gamma_{ki}^\eps. 
$$
Then the Gauss-Bonnet  argument (as in Proposition \ref{thick part estimate}) shows that 
$$
    \sum l_\hyp ( \gamma_k^\thick ) = O(\area (P) ) = O(|\chi | ).
$$
Putting this together with (\ref{geod and annuli}) and (\ref{rectangle  and geod}), we obtain the
desired. 
\end{proof}

% \bignote{``Marked Riemann surfaces'' are hidden here} 

\subsection{Marked Riemann surfaces with boundary}
\note{need?}

Two particular cases that we have considered above can be unified in
one frame. 

\sss{Boundary marking}
Assume now that some boundary components $J^k$ of $S$ are   
partitioned into several intervals, 
$\II^k= (\BI_1^k, \dots, \BI_{p_k}^k)$.
If some component $J^k$ is not
partitioned, we let $J^k\equiv I^k_1$  and  $p_k=1$.
% In this case we can still refer to $J^k$ as an ``interval''.
% We will also use one-index labeling $I_n$ if it does not lead to a confusion.  
We say that $S$ is endowed with  the boundary marking $\II=\bigsqcup I^k$.
We let $|\II| = \sum p_k$ be the {\it cardinality} of the marking.  
 
An {\it arc} $\alpha $ on a marked surface is a path connecting some
{\it non-adjacent }  marked 1-manifolds $I^k_m$ and $I^j_n$, %  (that could be circles),
up to proper homotopy.  Lifting $\alpha $ to the universal covering $\D$,
we obtain an arc $\Bsi$ connecting two ideal intervals $\BJ^k_n$ and
$\BJ^i_m$.
Then we define the weight of $\alpha$ as

$$
   \WW(\alpha)= \WW (\BJ^k_n, \BJ^j_m). 
$$

Arcs with positive weights form a weighted arc diagram $\AAA(S, \II)$
of the marked surface.


Let us also define {\it local weights} $\WW( J^k_n)$.
If $J^k_n$ is a circle then the definition is the same is in the
unmarked case. Otherwise we define it as the width of the recatnge in
the universal covering with horizontal sides $\BJ^k_n$ and the
complement in $\T$ of $\BJ^k_n$ together with two adjacent lifted intervals. 

The {\it total weight} of a marked surface is defined as
$$
   \BW(S, \II) = \sum_\II \WW( J^k_n ).
$$ 


\sss{Thin-thick decomposition}

\begin{thm}\label{thin-thick for marked surfaces}
    Let $(S, \II)$ be a marked  compact Riemann surface with bundary,
$\chi=\chi(S)$ be its Euler characteristic. Then
$$
  \BW(S, \II) - O( | \chi |  + |\II|) \leq 2\, \sum_{\AAA(S, \II) } \WW(\alpha ) \leq
  \BW(S, \II).
$$   
\end{thm}

\note{Use total weight notation}

\begin{proof}
  Combination of the above two arguments. 
\end{proof}

\sss{Proper domains}
% \sss{Polygons}
Let $S$ be  a compact Riemann surface $S$ with  boundary, and 
let $D$ be a   % simply connected
domain in $S$ (defined up to homotopy) bounded (rel $S$) by several
ideal intervals  $\tl J_n\subset  J_n $ and several arcs $G_j$ (called {\it edges}).
% We will call $D$  a {\it proper polygon}.
 We  call $D$ a {\it proper domain}.)
 
% Let $\BD$ be a lift of $D$ to the universal covering $\D$.
% Its ``horizontal'' sides  are contained in some
% intervals $\BJ^n\subset \T$ that cover $J_n$, which determines a
% partial marking of $\T$. Gaps $\BG_k$ of this marking correspond to
% the edges $G_k$ of $D$.   

Let $\BD$ be 
 the covering of $S$ corresponding to  $D$. It is  a partially marked
 Riemann surface $(\BD, \II)$  with  marked intervals  $\BJ^k_n$
 covering the $J_n$,
 together with other components of $\di D$.
%%%  and gaps $\BG_j$ covering (with degree $1$)  the $G_j$. 

\bignote{Pay attention to  the outer boundary}
 
This determines local weights $\WW(\tl J_n) \equiv \WW(J_n) : = \WW (\BJ_n)$,
{\it  fluxes} $\WW^\perp (G_j):= \WW(\BG_j)$, and the total weight of $D$:
$$
   \BW(D)  : = \sum \WW (J_n) + \sum \WW^\perp (G_j).  
$$

% We can also define $\BW^{gap} (D)$ by .....

\msk
We can also consider arcs $\alpha$ in $D$
connecting various pairs of intervals $\tl J_n$ and $\tl J_m$.
Letting   $\WW(\alpha):= \WW(\ba)$, 
where $\ba$ is the  lift of $\alpha$ to $\BD$, we obtain the weighted arc
diagram $\AAA_D$.

Along with arcs, we can consider {\it segments} $\si$ connecting two
edges $G_j $ and $G_i$.  Each segment is lifted to an arc $\Bsi$ in $\BD$
connecting the corresponding  gaps  $\BG_j$ and  $\BG_i$. 
Letting $\WW(\si) := \WW(\Bsi)$, we obtain  the  {\it weighted segment diagram}
$\SSS_D$. 

We  can also consider {\it arc-segments} in $D$ connecting intervals
$\tl J_n$ to  gaps $G_j$, and to define their weights.
They form the  {\it weighted  arc-segment diagram} $\AS_D$. 
Altogether,
arcs, segments and arc-segments form the  {\it weighted  arc-segment diagram}
$\SAS_D$. 

\ssk
Summing  up the  weights of arcs, segments, or arc-segments,
we obtain the {\em total weights} of the corresponding diagrams,
$\BW(\AAA_D)$, $\BW (\SSS_D)$, and $\BW(\AS_D)$.
{We will use the same notation, $\BW(\QQ)$, for any famly of
 arc-segments.)  
Theorem \ref{main estimate}
% \ref{thin-thick for marked surfaces}
implies:

 \begin{cor}\label{thin-thick for domains}
  Let  $D\subset S$ be  a proper $p$-gon as 
   Then  
 $$
      \BW(D) - O( p)  \leq   2\, \sum_{\AS_D}  \WW(\si)    \leq  \BW(D). 
 $$
 \end{cor}


\comm{*****
In fact, for any arc $G$ in $S$  we can define the flux 
$\WW^\perp(G)$ through $G$ as follows. 
% by  letting $D$ be the surface $S$ cut along $G$.
Let $G$ land on boundary components $J_n$ and $J_m$.
Then it can be lifted to an arc $\BG$ in $\D$ landing on $\BJ_n$ and
$\BJ_m$.
This determines an ideal  quadrilateral with {\it vertical} sides $\BJ_n$
and $\BJ_m$. Its width is $\WW^\perp(G)$.  

Similarly, given two arcs $G_i$ and $G_j$ 
(that could be equal: $G_i=G_j$)
 connected by  a segment
$\si$ in $S$, we can define the weight $\WW(\si)$  as follows. 
% by considering a
% domain $D$ obtainded by cutting $S$ along $G_i$ and $G_j$. 
Let $G_i$ land on $J_n$ and $J_m$, while $G_j$ land
on $J_s$ and $J_t$. Then we can lift the whole configuration to $\bar\D$
marking four  $\BJ$-intervals on $\T$. The segment $\Bsi$ connects two
gaps in this configuration that should be viewed as the horizontal
sides of a quadrilaterlal. Then $\WW(\si)$ is the width of this quadrilateral.


Of course, we can deine the weight of a general arc-segment in $S$ as well. 

******* } 

\subsection {Transformation rules}

By passing to the Universal Covering (or to the annulus covering), 
transformation rules of  Lemma \ref{transform rules on D} readily
generalize from the disk to arbitrary marked Riemann surfaces: 

\begin{lem}\label{generaltransform rules}
Let $i: S' \ra S$ be a  holomorphic map between
two compact  Riemann surfaces with boundary.
Then 

\ssk\nin $ \mathrm{(i)} $  If a boundary component $J' $ of $S'$ is mapped
with degree $d$ to a boundary component $J $ of $S$ then 
$ \WW(J')\leq d\cdot \WW(J)$;

\ssk\nin $ \mathrm{(ii)} $  If a  gap  % segment
$G ' $ of $S'$ is mapped
to a   gap  % segment
$G $ of $S$ then $\WW^\perp (G')\geq \WW^\perp (G)- 2 $;

\bignote{Used to be a ``segment'' instead of a ``gap''.
  Also: Notation``$\si$'' is used for segments} 

\ssk\nin $\mathrm { (iii) } $
 If an arc $a ' $ of $S'$ is mapped
to an arc  $\alpha $ of $S$ then $\WW (\alpha')\leq \WW (\alpha)$;

\ssk 
In case of a covering map, all the above inequalities become
equalities (without  subtracting ``$2$'' in $\mathrm { (ii) } $). 
\end{lem}

% \begin{cor}
%  Under the circumstance of the above lemma, let $D'\subset S'$ and
%  $D\subset S$ be  proper domains such that $i(D')$ is properly
%  isotopic to $D$. Then 
% $$
%   \BW^{gap} (D') \geq \BW^{gap}(D) - O(p+|\chi|).
% $$ 
% \end{cor}

\bignote{Extension to arc diagrams is hidden here} 

\comm{*****
We can also naturally define the flux $\WW^\perp (X)$  through a weighted arc diagram 
$X= \sum x_i G_i$, namely
$$
         \WW^\perp (X) =\sum x_i \WW^\perp (G_i).
$$
By linearity, property (ii) of Lemma \ref{generaltransform rules}
extends to this more general  situation:

\begin{cor}\label{flux through arc diag}
  Under the circumstances of the above lemma,
if  a weighted  arc diagram $X ' $ of $S'$ is mapped by $i$
to a weighted  arc diagram   $X $ of $S$ then 
$$
  \WW^\perp (X')\geq \WW^\perp (X)-O(|X'|),
$$
with the equality (with no error term) in case of the covering map.
\end{cor}
*****************}


Finally, let us also consider the case of concatenation of several
arcs. By Lemma \ref{generaltransform rules} \note{specify} 
and  the Parallel Law,we have: 


\begin{lem}\label{transform rule for concatenation}
Let $i: S' \ra S$ be a  holomorphic map between
two compact  Riemann surfaces with boundary.
Let $G_1', \dots , G_n'$ be a family of arcs on $S'$ 
whose concatenation through some imporper boundary components $J_k'$ 
forms an arc $G$ in $S$. Then
$$
   \sum \WW^\perp (G_k') + \sum \WW( J_k') \geq \WW^\perp(G) - O(n). 
$$
\end{lem} 


\sss{Key estimate} 

Let $S$ and $S'$ be two Riemann surfaces, and let $D\subset S$ and
$D'\subset S'$ be two proper domains, endowed with their
segment diagrams, $\SSS_D$ and $\SSS_{D'}$.
% immersed foliations $\FF$ and $\FF'$. 
%Note that the domains $D$ and $D'$ have
%the same cardinality $p$ of their markings $\II$ and $\II'$, 
%as well as they have the same Euler characteristic.

Let $i: S' \ra S $ be a holomorphic map such that $i(D')$ is properly
  isotopic to $D$.  Given any  segment $\sigma \in \SSS_D$, the pullback
  $i^*(\si)$ is a   concatenation of several  segments $\si_k'\in \SSS_{D'}$.  \note{picture?} 
  We say
  that $\si$ {\it breaks} under the pullback if $i^*(\si)$ comprises more
  than one segment. Let $\SSS_D^\br$ stand for the set of segments that break.     
 % Following our notational style, $\BW (\SSS_D^\br)$ means their total
%  weight. 
  

% Given any  leaf $L$ of $\FF$, the pullback
%  $i^*(L)$ is a   concatenation of several  segments in $D'$. We say
%  that $L$ {\it breakes} under the pullback if $i^*(L)$ comprises more
%  than one segment. Let $\FF^\br$ stand for the set of leaves that break.     
  
%As usual, the weights of all path families are measured in the
% Universal Covering. 


\begin{lem}\label{key estimate} 
  Under the above circumstances,  we have:
$$
     \BW (\SSS_D^\br) \leq \BW (\SSS_{D'}) - \BW (\SSS_D)  + O(p+|\chi|). 
$$
\end{lem}

\begin{proof}
  By Maximality of $\SSS_{D'} $ we have:
$$
\BW(\SSS_{D'}) \geq \BW  ( i^*(\SSS_D) ) - O(p+|\chi|)
  \geq \BW (\SSS_D)  + \BW(\SSS_D^\br) - O(p+|\chi|).
$$
The last estimate follows from the fact that    $\WW(\si') \geq \WW(\si)$
for any lift $ \si_k'\in \SSS_{D'}  $ of  a segment $\si\in \SSS_D$. Since
broken segments  admit at least two lifts, their weights are counted
at least twice in $i^* (\SSS_D) $, so
$ \WW(i^*(\SSS_D)) \geq \WW(\SSS_D)+\WW(\SSS_D^\br)$.  
\end{proof}


\comm{*****
\begin{lem}\label{key estimate} 
  Under the above circumstances,  we have:
$$
     \WW(\FF^\br) \leq \BW(D', \II') - \BW (D, \II)  + O(p+|\chi|). 
$$
\end{lem}

\begin{proof}
  By Maximality of $\FF'$, we have:
$$
   \BW(D', \II') \geq \WW( i^*(\FF) ) - O(p+|\chi|) \geq \BW (D, \II)
   + \WW(\FF^\br)
   - O(p+|\chi|).
$$
The last estimate follows from the fact that    $\WW(\si') \geq \WW(\si)$
for any lift $ \si' $ of  a segment $\si$ in $D$. Since
broken segments $\si$ admit at least two lifts, their weights counted
at least twice in $i^* (\FF)$, so
$ \WW(i^*(\FF)) \geq \WW(\FF)+\WW(\FF^\br)$.  
\end{proof}
************************} 

We can similarly define arcs and arc-segments in $D$ that break,
$\AAA_D^\br$ and $\AS_D^\br$.
% and their total weights,
% $\BW (\AAA_D^\br)$ and $\BW (\AAA_D^\br)$ .
Let us say that an arc in $D$  is {\em short} if it connects 
 two consecutive intervals, and {\em long} otherwise.
 Let $\AAA_D^\sh$ and $\AAA_D^\lo$ stand respectively for the
 familes of those arcs, and let $\AAA_D^\lobr$ stand for the
 family of long arcs that break.


 \begin{lem}\label{broken arc-segments} 
   Under the above circumstances, we have:

 %  \ssk\nin {\rm (i)}
$$
\BW(\AS_D^\br) \leq \BW (\SSS_{D'} ) \quad {\mathrm {and}} \quad
\BW(\AAA_D^\lobr) \leq \BW (\SSS_{D'} ) + \BW (\AS_{D'} )  .
$$
   
 \end{lem}

 \begin{proof}
The pullback $i^*(\alpha) $ of any arc-segment  $\alpha\in \AS_D^\br$
 consists of at least two pieces, one of which must be a
 segment in $D'$. Moreover, the weight of this segment is at least
 $\WW(\alpha)$. Hence
 $$
 \BW(\SSS_{D'} ) 
%  \sum_{\alpha\in \AS_D^\br} \WW  (i^*(\alpha) )
 \geq \sum_{\alpha\in \AS_D^\br} \WW (\alpha) = \BW(\AS^\br).  
 $$

 \ssk 
Similarly,    the pullback $i^*(\alpha)$  of any  long arc
$\alpha\in \AAA^\lobr $ that breaks consists of

\ssk\nin $\bullet$ either two pieces, one which must be an arc-segment
in $D'$;

\ssk\nin $\bullet$
or at least three pieces, one of which (a ``middle'' one)  must be a
segment in $D'$.

\ssk The second desird estimate follows. 
\end{proof}   
 
 
\comm{****** 

\subsection{Closed curve diagram}

\sss{Definition}
Given a non-trivial and non-cuspidal  simple closed curve $\gamma$ on a Riemann surface $S$,
let  $\A_\gamma$ be  the covering annulus  corresponding to $\gamma$,,
and let  $\Bga$ be the closed geodesic homotopic to $\gamma$.
Then we let 
$$
    \WW(\gamma)= \mod \A_\gamma= \pi/  l_\hyp(\Bga).
$$

For $M$ sufficiently big (``4''), 
the $M$-{\it canonical curve diagram} $\CC$ is the family of simple
closed curves $\gamma$ with $\WW(\gamma) > M$. 

\begin{lem}\label{disjointness in CC}
  The curves is $\CC$ are pairwise disjoint.  
\end{lem} 

\sss{Transformation rules} 

\begin{lem}\label{tranform rules for CC}
  Let $i: S' \ra S$ be a holomorphic map between Riemann surfaces.
Let $\gamma\in \CC$ and $\gamma'\in \CC'$  be simple closed curves
such that $i(\gamma')$ is homologous to  $d\cdot[\gamma$. 
 Then $\WW(\gamma') \leq d^{-1} \cdot \WW(\gamma)$.
In the case when $i$ is a covering, the equality holds. 
\end{lem}

\begin{proof}
We can assume that $\gamma'$ is a closed geodesic.
By the Schwarz Lemma, 
$$
     l_\hyp(\gamma') \geq  l_\hyp (i(\gamma)) = d\cdot l_\hyp (\gamma),
$$
and the conclusion follows. 
\end{proof}

\sss{Flux vs weight}

\begin{lem} \label{flux vs weight}
  Let $\alpha $ be an arc, and let $\si$ be a segment with both ends on
  $\alpha$ that land on the opposite sides of $\alpha$ (the
  orientation preserving situation).  Let us consider a simple closed curve $\gamma$ obtained 
  by connecting the endpoints of $\si$ along $\alpha$. Then
$$
   \WW(\gamma) \geq \frac 12  \bar\WW (\si), 
$$  
as long as $\bar \WW(\si) > 2/\pi$.

\end{lem}

\begin{proof}
  Let $\alpha$ lenad on boundary components $I$ and $J$.
Then our configuration of arcs lifts to the universal covering $\D$ as
follows:

\begin{itemize}
\item $I$ lifts to two ideal arcs $\BI$ and $\BI'$; 
          $J$ lifts to two ideal arcs $\BJ$ and $\BJ'$

\item $\alpha$ lifts to two arcs: $\Bal$ connecting $\BI$ to $\BJ$,
  and $\Bal'$ connecting $\BI'$ to $\BJ'$. 

\item $\si$ lifts to an arc $\Bsi$ connecting $\Bal$ and $\Bal'$;

\item $\gamma$ lifts to aa axis $\Bga$ connecting the ideal gap $\BG$
  between $\BI$ and $\BJ$ to the ideal gap $\BG'$ between $\BI'$ and $\BJ'$
  (and separating $\BI$, $\BI'$ from $\BJ$, $\BJ'$.

\end{itemize} 

Let us uniformalize $\D$ be the horizontal strip of width $\pi$ so
that $\Bga$ becomes an axis of this strip. (We will keep the same  notatin 
for all the objects under consideration.) The deck transformation
$T_\gamma$
along this axis is a translation by some number $a$. Moreover,
$\WW(\gamma) = \pi/a$. 

The curves representing $\WW(\si)$ connect $\BG$ to $\BG'$,
so they go over the convex hulls  $[\BI, \BI']$ and $[\BJ, \BJ']$,
both of which have length greater than $a$. Hence these curves have
length at least $\min (\pi, a)$. 

Let us put the Euclidean metric on the union of two rectangles based upon the above
intervals
and going across the strip. Then the length of all curves with respect to this
metric   is bounded in the same way. The area of these metrics is
bounded by $2\pi a$. Hence   
$$
  \bar\WW(\si) \leq \frac {2\pi a} { \min (\pi, a)^2  } =\max \left(\frac{2\pi}
  a, \frac 2\pi \right) = 2 \max ( \WW(\gamma), 1/\pi). 
$$
The conclusion follows. 
\end{proof}

******* } 

\comm{***** 

\sss{Flux through the ray diagram}

For an arc $\alpha$ on $U'\sm \KK'$, the inclusion $i_*(\alpha)$ is well
defined if

\ssk\nin $\bullet$
$\alpha$ is horizontal with both ends landing on some $K_n$; 

 \ssk\nin $\bullet$
or $\alpha$ is vertical with one end landing on some $K_n$; in which case it admits a unique extension
to a vertcal arc $ i_*(\alpha)$ on $U$ by means of the homotopy
retraction $U\ra U'$. 

\msk
We say that a WAD $X'$ on $U'$ is {\it invariant} under $F$  if :

\ssk \nin $\bullet$ 
The inclusions $i_*(\alpha)$ are well defined for all $\alpha\in
\supp X'$; 

\ssk \nin $\bullet$ $f_*\alpha$ is an arc for any $\alpha\in \supp X'$
  (in other wirds $\alpha$ and $-\alpha$ do not cross); 

\ssk \nin $\bullet$   $i_*(X')= f_*(X')$.
 

\ssk\nin
We let $X=f_*(X')=i_*(X')$ be the image WAD on $U$. 

\msk
For instance, consider the {\it ray diagram}  $\RR'$ on $U'\sm \KK'$ supported on all external rays
landing on the little Julia sets $K_n$, $n =0,\dots, p-1$, with weights
$$
   r_n' =  (\# \ \mbox{ of external rays landing on $K_n )^{-1}$. }
$$

\begin{lem}\label{flux through inv WAD}
    Let $X'=\RR'$ be the  invariant ray daigram on 
$U'\sm f^{-1}    (\KK)$. 
Then
$$
  \WW^\perp (Y')\geq \WW^\perp (X') - O( p (1 +  \WW^\ver) ) 
% - o(p\, \WW_\loc) 
$$ 
for any   WAD $Y'$  such that  $i_* (Y') = i_* (X')$.
\end{lem}

\begin{proof}
Let  us decompose $i$ as $j\circ e$ where
$$
  e : U'\sm (\KK\cup \KK') \ra U'\sm \KK,\quad j: U' \sm \KK \ra U\sm \KK.
$$
We let $Z = e_*(X')$, so $j_*(Z) = X$.   Then
%
\begin{equation}\label{restricted X}
   \WW^\perp (X) \leq \WW^\perp (Z) + O(p (1 +  \WW^\ver) ). 
\end{equation}
Indeed,  for any proper path $\gamma$  in $U\sm \KK $  crossing $\alpha \in \supp X$, 
a lift $j^*(\gamma)$ that begins on a little Julia setis either
crosses $ j^*\alpha\in Z$  or is vertical
 (and the vertical weights
get summed up over the whole diagram $X$ containing $O( p) $ arcs). 

\comm{******
If it is vertical, three situations can occur:

\ssk\nin (i) $\gamma$ is well localized in the translation region
(meaning that its ``dives" into the Translation region withon a
bounded number of blocks). The total weight of such  

\ssk \nin (ii) $\gamma$ dives into the Translation region in a
non-localized way. The total weight of such arcs is
$O(p+\WW_{periph})$. 

\ssk \nin (iii) $\gamma$ exits the Translation region.
The total weight of such arcs is
$O(\WW_{periph})$. 

% \bignote{What to do with the peripheral error ?}
*********}

Since $i_*(Y') =X$, we have $e_*(Y') = Z$, and Corollary 
\ref{flux through arc diag} yields: 
$$
   \WW^\perp (Y') \geq \WW^\perp (Z) - O(p).
$$
Putting it together with (\ref{restricted X}), we obtain: 
$$
      \WW^\perp (Y') \geq \WW^\perp (X) - O(p(1+  \WW^\ver)) 
% - o(...).
$$
But since $f: U' \sm (f^{-1}\KK)\ra U\sm \KK$ is a covering map and
$f_* (X')= X$, 
we also have:
$$\WW^\perp (X') = \WW^\perp (X).$$ 
\end{proof}
   
\begin{rem}\label{refinement}
The above estimate comes through if we replace the Riemann surfaces
$S=U\sm \KK$ and $S'= U'\sm f^{-1} \KK$ with 
$$
    S^n  = U^n\sm f^{-n} \KK. \quad \mathrm{and} \quad   S^{n+1}   
$$
coming together
with the  covering-immersion pair $F =(f ,i) : S^{n+1} \ra S^n$, 
and two   arc diagrams,  $X^{n+1} =\RR^{n+1} $ and $Y^{n+1}$,  on $S^{n+1}$ 
such that $X^{n+1} $ is $F_n$-invariant and  $i_ * (Y^{n+1} ) = i_*(X^{n+1})$.   
The error term should be interpreted as the vertical weight on $S^{n+1}$. 
When we want to emphasize  where the weight is calculated, 
we use notation like $\WW(X; S)$. 
\end{rem}

Next, we will sharpen the error term in  the previous lemma. 


\begin{lem}\label{flux through inv WAD-improved}
Fix an arbitrary small $\eps>0$ and a big $N$. 
Under the circumstances of the previous lemma
(adjusted as in the above Remark), 
either $W^\ver \geq N \WW^\loc$
or else
$$
  \WW^\perp (Y^n;\, S^n  )\geq 
\WW^\perp (X^n; \,    S^n) - \eps\, p\, \WW^\loc  
$$ 
for some $n = O( N/\eps) $. 
% (Here $U^n$ is the $n$ times restricted domain).
\end{lem}  

\begin{proof}
We will pursue the following strategy. 
  If for some restriction to  $S^n$, 
the newborn vertical saisfies $W^\ver (S^n) < \eps \WW^\loc$, 
then we obtain the desired estimate. Otherwise, we have the
definite (in terms of $\WW^\loc$) vertical weight created under each restriction. 
Under $\asymp N/\eps$ steps, we will create at least $N\WW^\loc$ of
the vertical weight.
Then we will push it forward by a baby version of the Covering Lemma. 


Let  
$$
\Si^n : = U^n \sm f^{-(n-1) } \KK.
$$

%For $n\geq N$, let  
%$$
%S^n  : = U^n\sm f^{-N} \KK, \quad
%(S^{n+1})' : = U^{n+1} \sm f^{-(N+1) } \KK.
%$$

Let us consider the covering annulus $\A^n$ of $S^n$ corresponding to the loop
$\di U^n $, and let $\BBB^n $ be the similar annuli for the $\Si^n $.

Let  $\Pi^n$ be the universal coverings of the $\A^n$  realized as strips
(with the top $T^n  =\R+ h_n \pi $  covering the outer boundary $\di U^n$ and the
bottom $\R$).  We  select the heights $h_n$ so that $\A^n= \Pi^n/\Z$. 
Similarly, $P^n$ are the univesal coverins of the $\BB^n $. 

We mark external rays  $\hat \RR^n(i)$  on the annuli
and the strips within one fundamental region
(using the same notation in $\A^n$ and $\Pi^n$). 
They select for us  initervals $\hat K^n(i)$ on the bottom $\R$ corresponding to the little Julia
sets. By definition, the flux $\WW^\perp$ through $\RR^n(i)$ is
equal to the height of 
% width of the vertical  foliation $\FF^n$ in $\Pi^n$  through $\hat\RR^n(i)$ in 
the quadralateral $\Pi^n(i)$ whose horizontal sides are
the top of $\Pi^n$  and the interval $\hat K^n(i)$. 

\comm{*****
Let $I^n$ be the interval on the bottom $\R$ of $\Pi^n$ which is the
convex hull of the  the intervals $\hat K^n ( q-(n-1) ) $ and 
$\hat K^n (q- (n-3))$. 
Let $w^n $ be the  the width of the path family going from $\R$ to
$\R$ skipping over $I^n $
(i.e., the height 
 of the rectangle $\Pi^n$ with
horizontal sides $I^n$ and $T^n$). 
We claim that the $w^n$ are small.
% in terms of $N$ only. 

To see it, let us also consider the
% annulus $\A$ covering $U\sm \KK$, its universal covering $\Pi$.
% Then we have  the 
covering map $f_n :  \A^n \ra \A$ of degree $2^N$ and its linear
lift $\hat f_n :  \Pi_N\ra \Pi$. 
Under this map,  the interval $\hat f_n  ( I^n) $ 
covers (essentially with degree one) the whole bottom of $\A$, which implies that $w_n$ is bounded
by the height of  $\A$ plus $O(1)$ 
% the rectangle $\Pi$ with horizontal sides $1$ and $T$. 
Since this annulus is bad, we are done. \note{comment?} 
*****}


Let $k=[1/\eps]$, and let $n\leq kN$. 
Let us consider the foliation $\FF^n (i)$ representing the flux
through $\RR^n (i)$ in $\Pi^n $. 
Let us pull it back by $i^*$ to $\Pi^{n+1}$.  
A leaf $L$ can either be fully lifted to $\Pi^{n+1}$ or else it can be broken.
 It the former case,  it contributes to the flux through $\RR^{n+1}(i)$.
It the latter case, it has at least  two vertical lifts that begin on
the opposite side of $K^n(i)$.  
% If some leaf $L$ of $\FF^n$ fully lifts to $\Pi^{n+1}$ then all the leaves
%below it do the same. It follows that the family of ``broken leaves''
%$i^*(L)$ form a foliation near the top. 

So, the weight of the broken leaves is bounded by the vertical weight
of $A^{n+1}$.  If the total of this flux  is smaller then $\eps p  \, \WW^\loc$ then 
% $$
%    \sum_i     \WW(\FF^N(i) ) -  \WW(\FF^n(i) )  < \eps (n-N) p \, \WW^\loc.
%$$
%Hence for some $k\in (N,n]$ we would have
%$$
%     \sum_i  \WW(\FF^{k-1}  (i) ) -  \WW(\FF^k(i) )  < \eps p \, \WW^\loc,
%$$
we obtain  the desired error estimate. 

Otherwise, for all  $n\leq k N$, the weight of broken leaves of
$\FF^n(i)$ (for any $i$ in the ``middle''  the translation region) is at
least $\eps\WW_\loc$. 
 % $$
%  \sum_i     \WW(\FF^N(i) ) -  \WW(\FF^n(i) )  \geq  \eps p (n-N) \, \WW^\loc,
%$$
%so for some $i$. 
%$$
%    \WW(\FF^N(i) ) -  \WW(\FF^n(i) )  \geq  \eps  (n-N) \, \WW^\loc,
%$$
After $kN$ restrictions we create at least $N\WW_\loc$ of the vertical
weight.  

If it is localized outside the interval $I^n_{kN}$ in $\A^{kN}$ contianing the
last $kN$ sectors, 
%As we know, only small amount of $\FF^n(i) $  can skip over $I^n$,
%so most of this flux lands well inside the trnalation region (at least
%on its forward side). Assume good part of it actually lands outside
%$I^{kN}$.  
then there is  a standard rectangle $P$ in the translation
region but outside $I^n_{kN}$ whose width is at least as big  (up to a bounded additive
error).  Under the map 
$\hat f_{kN}: \Pi^{kN} \ra \Pi$ composed with the projection $\Pi\ra \A$, 
this rectangle is mapped with degree one to $\A$.
Hence  $A$ has at least as big width.  
%
%By the above consideration, $\WW (\A^N)$ is at least $\asymp N\WW^\loc$.
%Moreover, a definite part of the newly created vertical weight is
%localized on a vertical rectangle $\Pi$ in $\A^N$ such that $f^N$ is
%injective on its horizontal side (since it can be selected well inside
%the translation region). By the Covering Lemma (a baby version), 
%this weight can be pushed forward to $\A$.   

Otherwise we would create a big vertical weight in $\A^n$ landing on
$I^n_{kN}$.  Let us consider a disk $D$ containg  $kdN/2$ little Julia
sets (in the $kN/2$ consecutive sectors) aligned with the Hubbard
marking.  Let $\mathbf D$ be the corresponding psuedo-disk. 
The above vertical weight blocks the canonical arc diagram in $\mathbf
D$ to skip over several little Julia sets in the same translation
orbit from ``above''. 
 Neither an arccan be a ``serpant'' winding above these little Julia
 sets
since it would create a self-intersection. So, any arc can leave
$\mathbf D$ only through $\RR(0)$ or  by going ``under'' all little Julia sets,

In either case, its push-forward
% would  connect the next little Juli set to $\RR(1)$ and 
would  intersect itself. 
\end{proof}

\sss{Cobras}


Let us begin with a non-dynamical situation. 
Let us mark an (oriented)  arc $\alpha$  in $U\sm \KK$ connecting a
boundary component  $I$ to a component $J$.  
A  segment $\si$ in $(U, \alpha)$  is called a {\it cobra}
(attached to  $\alpha$) if it  begins and  ends  on $\alpha$ and is
homotopically trivial in $U\sm (I \cup J \cup \alpha)$ (where its endpoints
are allowed to slide along $\alpha$).  
It is oriented so that its end is ahead of the tail (on $\alpha$).  
 Its weight
$\WW(\si)$ is the flux through $\alpha$ along $\si$.  \note{defined?}


 Let us consider a twisted arc $\alpha^\si$ on $U\sm \KK$  obtained by adding to $\alpha$
 the  segment $\si$, while erasing the corresponding piece of $\alpha$.  \note{picture}
The following lemma shows that this operation leads to a loss of the
flux (as long as the weight of the cobra is sufficiently big). 

\begin{lem}\label{loss of cobra}
If $\WW(\si) > 0 $ then 
$$
 \WW^\perp (\alpha^\si)  \leq  \WW^\perp (\alpha) - \WW( \si ) + O(1).
$$
\end{lem}

\begin{proof}
   Let us lif the segment $\alpha$ to a segment $\ba$ connecting some
   ideal intervals $\BI$ and $\BJ$ (lifts of of the corresponding
   boundary components). 
   Let $\Bsi$ be the lift of the cobra that begins on $\ba$. It
   connects $\ba$ to another lift $\ba_1$ of $\alpha$.  
  The latter begins on some ideal interval  $\BI_1$ (another lift of $I$)  and ends on some
  $\BJ_1$.

Assume for definiteness that the cobra $\si$ is negatively oriented
with respect to $\alpha$ (i.e.,  we turn right from $\alpha$ to
$\si$). Then the ideal intervals under consideration appear on the
circle $\T$ in the following cyclic order: $(\BI, \BJ_1, \BI_1, \BJ
)$. Let us consider the canonical arc-segment diagram of this partially marked
disk. By our assumption, it contains the arc $\Bsi$ connecting the gaps 
$(\BJ, \BI)$ and $(\BJ_1, \BI_1)$. Besides, it may contain four segments $\Bga(\BI), \dots,
\Bga(\BJ)$ connecting the gaps attached to the corresponding intervals
$\BI$ etc., and  some arc-segments begining on these
intervals (altogether there are at most five arc-segments with
positive weight).   

 Let us first assume that all the latter arc-segments have zero weight.
 Then we have:
$$
   \WW^\perp(\ba) = \WW(\Bsi) + \WW(\Bga(\BI)) + \WW(\Bga(\BJ)) + O(1),  
$$    
$$
   \WW^\perp(\ba^\si) = \WW(\Bga(\BI))  + \WW(\Bga(\BJ_1)) +O(1). 
$$      
So, it is enough to show that
\begin{equation}\label{desired estimate} 
   \WW(\Bga(\BJ_1)) \leq \WW(\Bga(\BJ) ).
\end{equation} 

Let $T$ be the deck transformation moving $\ba_1$ to $\ba$. 
It moves the configuration $\CC:=(\BI, \BJ_1, \BI_1, \BJ)$ to 
$T(\CC)= (\BJ, \BI,  T(\BJ),   T(\BI) ) $ (up to cyclic permutation).
Let $L$ be the gap  between $T(\BI)$ and $\BJ$, and let $\Bde$ be the
segment of $T(\CC)$  connecting $L$ to  the gap $(\BJ, \BI)$. 
Then 
$$
  \WW(\Bga(\BJ_1)) = \WW(\Bde).  
$$
Moreover,  the  fixed points $q^u$ and $q^s$ of $T$ lie in the gap $(\BI_1, \BJ)$,
(with $q^s$ ahead of $q^u$). \note{to be justified}
It follows that $T(\BI)$ lies in  between $q^s$ and $\BJ$,
and thus  $L\subset (\BI_1, \BJ)$.
Hence $\WW(\Bde) \leq \WW(\Bga(\BJ) )$, and 
the desired estimate (\ref{desired estimate}) follows. 

\msk In the general case, there could also be the arc-segment $\Bom$ connecting
$\BI$ to the gap $(\BJ_1, \BI_1)$ that intersects $\Bsi$ but not
$\ba$. 
But then $T(\Bom)$ goes from $T(\BI)$ to $(\BJ, \BI)$  contributing
the same flux through $\ba$.
\end{proof}

 Let us now consider  a family $\SSS$ of disjoint cobras $\si_i$ with
 the same (say, negaive) orientation attached to
 some arc $\alpha$. Let $D_i$ be the Jordan disk bounded by $\si_i$ and
 the segment of $\alpha$ sharing the endpoints with $\si_i$.   
Any pair of these disks are  either nested or disjoint,
so we can talk about  innermost and outermost cobras. 
The outemost cobras form a  family with disjoint $D_i$.
Let us define the twisted arc $\alpha^\SSS$ 
by twisting $\alpha$ simultaneously  along the outermost cobras. 

We define $\WW(S)$ as the total weigh of the cobras in $\SSS$. 

\begin{cor}\label{loss of several cobras}
Let  $\SSS$ be a family of negative cobras of positive weight attached to some arc $\alpha$. Then 
$$
 \WW^\perp (\alpha^\SSS)  \leq  \WW^\perp (\alpha) - \WW( \SSS ) + O(p).
$$
\end{cor}

\begin{proof}
  The twisted arc $\alpha^\SSS$ can also be obtained by twisting
  $\alpha$ consecutively along the innermost cobras. In this way we
  will make use of all the cobras of $\SSS$, 
loosing each time the weight of the cobra, up to a bounded additive
constant
(Lemma \ref{loss of cobra}). Since there exists at most $O(p)$
disjoint cobras, the conclusion follows.  
\end{proof}

More generally, we can consider a disjoint family $\SSS$  of  cobras
with the same orientation attached to
a WAD $X$, and twist $X$ along this family.  
It immediately follows that 
\begin{equation}\label{cobras attached to diag}
    \WW^\perp(X^\SSS) \leq \WW^\perp (X) - \WW(\SSS) + O(p), 
\end{equation}
where $\WW(\SSS)$ should be computed with the corresponding weights. 

Let us now pass to a dynamical setting. 
Let $X$ be an invariant WAD in $U'\sm (\KK \cup \KK')$.
We say that the cobra $\si$  attched to $X$  is {\it tame} if $i_*(\si)$ is a trivial
segment on $U'\sm \KK$
(so, $\si$  surrounds only isles of $\KK'$). 

\begin{lem}\label{tame cobras}
Fix a big $N$.
Let $X$ be an invariant WAD, 
% with $\WW^\perp (X') \geq C\, |X'|$. Then
and let $ \SSS $  be a family  of disjoint  tame cobras $\si_i$
attached to $X$.  
If $\WW^\loc$ is sufficiently big then either $\WW^\ver\geq N \WW^\loc$ or else
$$
      \WW (\SSS)   = o ( p \WW^\ver ).   
$$
\end{lem} 

\begin{proof}
We can consider separately negatively oriented and positively oriented
families. In either case,  the twisted WAD $X^\SSS$  is  well defined. 
By (\ref{cobras attached to diag}),  
 we loose the total (weighted) flux of these cobras.
%$$
% \WW^\perp (Y')  = \WW^\perp (X') - \WW({\SSS} )  
% \leq (1-\eps) \WW(X')
%$$
% (assuming the desired estimate fails). 
On the other hand, since the cobras are tame, we have $i_*(X^\SSS) =
  i_*(X)$, so by Lemma \ref{flux through inv WAD-improved}
$$
     \WW^\perp (X^\SSS) \geq \WW^\perp (X) -  o ( p  \WW^\loc).
$$
and the conclusion follows. 
%Comparing this with the previous formula, we obtain the desired.
% $$   
%     \WW^\perp (X')  \leq \frac A\eps\,  |X'|, 
%$$
% which contradicts our assumption as long as $C> A/\eps$. 
\end{proof}


\comm{****
\sss{Generalized cobras} 
We will need  more general creatures:
a {\it generalized cobra} (attached to $X'$)  is a segment $\si'$ in
$U'$  with endpoints on two adjacent edges $\RR_1$ and $\RR_2$  of
$X'$ (so, both $\RR_i$ land on the same little Julia set $J_n$).
A generalized cobra is tame if in the region bounded by $\si \cup
\RR_1\cup \RR_2 \cup J_n$ (where $J_n$ is understood as a piece of
the ideal boundary  of $U'$) 
there are no other little Julia sets $J_k$.   
Lemma~\ref{tame cobras} has the following generalized version': 

\begin{lem}\label{tame generalized cobras}
Let $X'$ be an invariant WAD, 
% with $\WW^\perp (X') \geq C\, |X'|$. Then
and let $ \SSS $  be a family  of disjoint  tame generalized cobras $\si_i$
attached to $X$.  Then
$$
      \WW (\SSS)   = O( |X'| + p\, \WW^\ver + |\SSS|\, \WW_\loc ) .  
$$
\end{lem} 

\begin{proof}
Proceed as in  Lemma~\ref{tame cobras}: 
 consider a WAD $Y'$ on $U'$  obtained by adding to $X'$
   segments $\si_i'$, while erasing the corresponding pieces of $X'$.  
   Then:
$$
 \WW^\perp (Y')  = \WW^\perp (X') - \WW({\SSS} )  
% \leq (1-\eps) \WW(X')
$$
% (assuming the desired estimate fails). 

  On the other hand, since the cobras are tame, we have $i_*(Y') =
  i_*(X')$, so by Lemma \ref{flux through inv WAD}
$$
     \WW^\perp (Y') \geq \WW^\perp (X') -  O( |X'|+ p\, \WW^\ver).  
$$
Comparing this with the previous formula, we obtain the desired.
% $$   
%     \WW^\perp (X')  \leq \frac A\eps\,  |X'|, 
%$$
% which contradicts our assumption as long as $C> A/\eps$. 
\end{proof}

****}



\sss{Cobras in the Elephant eye}

Assume now that we have a multi-elephant combinatorics with $d$ eyes
(so the period is equal to $pd+1$. 
For the WAD $\RR'$ of external rays,
we let $\RR' (n)$ be the union of rays landing on $J_n$, and $\RR_k' (n)$
be the corresponding  individual rays.
Similarly,  for the diagram $\RR=i_*(\RR')$,
we use notation $\RR(n)$ and $\RR_k (n)$. 

Note that  $\RR (0)$ comprises two rays, so together with $J_0$ they
divide $U$  into two half-disk. Let $H_+$ be the one that does
not contain  the critical value $c$. For our combinatorics, it contains only
bounded number of the little Julia sets. Accordingly, $H_-$ is the
other half-disk. 

% We will use notation $\RR''(n)$ etc. for the rays in $U''\sm f^{-2} (\KK)$.

\begin{lem}\label{many tame cobras} 
 Any cobra $\si$ in $U \sm f^{-1} (\KK)$ attached to  $\RR(0)$ produces a tame cobra
 of the same weight in $U' \sm f^{-2} (\KK)$ attached to $\RR'(d\, p)$. 
\end{lem}

\begin{proof}
   Since the cobras in $U\sm f^{-1}(\KK)$ attached to $\RR'(0)$ form a $0$-symmetric 
family, we can assume without loss of generality that $\si$ is
contained in $H_+'$. 

    We have: $f(H_+')\supset H_+$,  % \note{for genuine ql maps}
and all little Julia sets $J_{pk}$ in $H_+'$ except $J_{pd}$ are mapped
to $H_-$. It follows that the  lift of $\si$ 
 by $f:  H_+' \ra U $  is a  tame cobra in $H_+'\sm f^{-2}(\KK)$ attached to
$\RR'(d\, p)$. 
\end{proof}

%  Comparing this with Lemma \ref{tame cobras}, we conclude: 

\begin{lem}
   If the total weight of cobras in $U\sm \KK$ attached to $\RR(0)$ 
is at least  $p \WW_\loc$ (in order), we are done. 
\end{lem}

\begin{proof}
%       We can assume that $\WW^\ver < 2 \WW_\loc$. 
%
% Reducing the cobras in question to  $U\sm f^{-1} (\KK)$ we loose
%$\leq p \WW^\ver$ of their weight. What is left is still much bigger
%than $p\WW_\loc$. 
    Lemma \ref{many tame cobras} produces for us a family of tame
  cobras in $U'\sm f^{-1}(\KK)$ of the total weight $p\WW^\loc$,
  \note{be careful}
and Lemma \ref{tame cobras} leads to a contradiction. 
\end{proof}


%\begin{rem}
%   How strong is the assumption that the total flux through $\RR(0)$ 
%is much bigger than $ p\WW_\loc $ ? 
%If it is violated then most leaves of the canonical foliation 
%intersect $\RR(0)$  bounded number of times. 
%I thnk this case can be reduced  the case of the CAD localized in the
%translation region, but have not yet looked at it closely.   
%\end{rem}

%{\it Question.}  
%How to rule out the situation when most of the flux through $\RR(0)$ 
%is comprised of curves that keep rotating around $0$ in the same
%direction ? 

\comm{****
\begin{lem}\label{many cobras}
Assume the canonical WAD is not localized in the translation region.    
At least $\asymp p\, \WW_\loc$  of the flux through $\RR(0)$ is supported on
 cobras in $H_+$.
\end{lem}

\begin{proof}
  In the non-localized case,  at least $\asymp p\, \WW_\loc$  of the flux through $\RR(0)$ 
is supported on the canonical WAD.   \note{justify}
Let us take an arc $\alpha$ of the
canonical WAD crossing $\RR(0)$. There are three possible cases:

\ssk \nin (i) It is trapped inside $H_+$;

\ssk \nin (ii) It returns to $\RR(0$ from $H_+$ crossing it through
the same ray as it entered;

\ssk \nin (iii)  
It returns to $\RR(0) $ from $H_+$ crossing it through
the diffenrent  ray.

\msk In case (i)  it lands on some  little Julia set in $H_+$
(possibly, on $J_0$), and the
total weihgt of such arcs is bounded by $(d+1)\, \WW_\loc$.  

\ssk
In case (ii) $\alpha$ bounds a cobra in $H_+$.

\ssk In case (iii), we have similar three subcases according to its behavior in $H_-$:

\ssk\nin (iii-a) After entering $H_-$ it is trapped over there; 

\ssk\nin (iii-b) It returns back to $\RR(0) $ from $H_- $ crossing it through
the same ray as it entered;

\ssk\nin (iii-c) It returns back to $\RR(0) $ from $H_- $ crossing it through
the different ray;

\ssk
Case (iii-a) leads to localization. \note{justify}

\ssk 
Case (iii-b) creates a cobra.

Finlly, in case (iii-c) we have again three subcases, etc. 

\msk In the end, we have: 

\ssk \nin 1)
 $\alpha$ either lands on some little Julia set in $H_+$:
totall weight of these arcs is still at most $(d+1) \, \WW_\loc$; 

\ssk \nin 2)  or it
creates a cobra based on $\RR(0)$ (which is what we want);

\ssk \nin 3) or it lands on some little Julia set in $H_-$ after a
bounded number of revolutions, which again leads to the localized case; 

\ssk \nin 4) or or it lands on some little Julia set in $H_-$ after an
unbounded number of revolutions.  
\end{proof}


\begin{lem}\label{many tame cobras} 
  Assume the canonical WAD is not localized in the translation region.   
Then  at least 1/2 of the flux through $\RR'(d\, p)$ is supported  on tame
    cobras. 
\end{lem}

\begin{proof}
    We have: $f(H_+)\supset H_+$, \note{for genuine ql maps}
and all little Julia sets $J_{pk}$ in $H_+$ except $J_{pd}$ are mapped
to $H_-$. Lifting by $f:  H_+'\ra U$  the cobras from  Lemma \ref{many cobras},  
 we obtain tame cobras in $H_+'$ attached to
$\RR(d\, p)$ that occupy at least 1/2 of the flux through $\RR'(d\, p)$. 
\end{proof}
****}

\comm{***********
\subsection{Rotation number} 

Let us define the {\it rotaion number}  $\rho(K)$ around a little  Julia set
$K=K_i$ as follows.   Take any ray $\RR$ landing at  $K$ oriented
from $K$ to $\di U$.   Consider the immersed foliation $\FF$
through $\RR$ oriented positively with repect to $\RR$. 
Then we have the (partially defined)   first return map $T: \RR\ra
\RR$. 
 For any $a\in \RR$ and $b= T(a)$, we let $c(a, b)= 1=
-c(b,a)$. If $T(a)$ is not well defined, we let $c(a,b)=0$ for all
$b\in \RR$.  
This function can be extended to a cocycle $c(a,b)$ for the
equivalence relation generated by $\FF$ on $\RR$. 
 Then we let
$$
   \rho(K_i) =  \limsup_{n\to \infty}   \frac 1n \int_{\RR} c(a,a_n) \, d\mu(a),
$$ 
where $a_n$ is the $n$th return of the leaf $L_a$ to $\RR$, 
and  $\mu$ is the transverse measure on $\FF$.

\begin{lem}
  The rotation number $\rho(K)$ does not depend on the choice of the ray
  $\RR$. 
\end{lem}
***************}

\subsection{Poincar\'e return map}

Given some arc $G$ for a map $f$, we can consider the return map $T=T_G$ to $G$ along the
immersed foliation $\FF^\perp$. 
More precisely, let us consider the covering corresponding to $U\sm
G$.  Its ideal boundary is marked with four ideal intervals, two of which, 
$\BG$ and $\BG^\#$, correspond to $G$.
% (two others correspond to the
% boundary components where $G$ lands). 
Let us consider the canonical foliation on this surface. 
It defines the  transit map from $\BG$ to $\BG\sqcup \BG^\#$,
which  can be viewed as the return map to $G$.  \note{makes sense?}
 It comes together wirh 
% the quasi-invariant transverse measure $\mu$ on $G$, and 
the sign function $\eps(x)\in \{0,\pm 1\}$ recording whether the return is well defined
(otherwise $\eps=0$) and if so, whether it is  orientation
preserving (in this case, it goes from $\BG$ to $\BG^\#$) or
reversing (then it  goes from $\BG$ to itself). 

\begin{lem}
  Assume we have a sequence of maps $f_n$ as above. Then the
  corresponding return maps converge almost everywhere to a
   measure preserving automorphism. 
\end{lem}

Here the ``measure'' $\mu$  is a weak limit of the flux measures
$\mu_n$ through $G$. 

If $G=\RR$ is an external ray then  $T$ is monotonic on $\{x: \ \eps (x)=1 \} $,
since the corresponding foliation (albeit immersed) \note{OK?} 
 lives on the plane. 


Given a sequence of maps $f_n$ with periods $p_n \to \infty$ 
and a ray  $\RR$ (independent of $n$ in an appropriate sense),
let us consider  the coresponding return maps $(T_n, \mu_n, \eps_n)$. 


\begin{lem}
  Assume 
$$  \WW^\ver (f_n) \leq 2 \WW_\loc (f_n), \quad p\cdot \WW_\loc (f_n) = o(\WW^\perp (\RR, f_n)); \quad
  \WW^\perp (\RR, f_n)\asymp  \locflux.
$$ 
Then the maps $(T_n, \mu_n, \eps_n)$ have a limit
$(T, \mu, \eps)$. Moreover, $\eps(x)\not= 0$ a.e., $T$ preserves $\mu$ and is monotonic on 
 $\{x: \ \eps(x)=1 \} $.
\end{lem}

\begin{proof} 
Let us orient  the ray $\RR$ from the little $J$ -set to $\di U$.
For any point $a\in \RR$,  orient the leaf $L_a^n$ of
$\FF_n^\perp$ positively with repsect to $\RR$.  

Let us consider the set $X_n\subset \RR$ of points that
  never return to $\RR$ in positive time. 
This set can return to another edge $G$ of our cell decomposition,
but it must be   wandering  under the return maps.  Hence the number
of returns is bounded by   
$$
       (1+\de)  \locflux) /\mu_n(X)
$$
(since the  trnasverse measure is almost preserved  under the
holonomies). 
It follows that there exists at most 
$$
   C p\cdot  \locflux  /\mu_n(X)
$$
different homtopy classes of arc-segments in this foliation (with some
absolute $C$). Since each of them has weight at most $\WW_\loc(f_n) $, we
obtain:
$$
  \mu_n(X)  \leq Cp\WW_\loc (f_n)\cdot  \locflux
  /\mu_n(X),
$$
and thus
$$
   \mu_n(X) \leq \sqrt{Cp\WW_\loc (f_n)\cdot 
      \locflux } = o(\WW^\perp(\RR, f_n )). 
$$
Thus, in the limit the return map is well defined almost everywhere. 
\end{proof}

\begin{lem}
  Under the assumptions of the previous lemma,  one of the following
  situations occur:

\ssk \nin $\mathrm{ (i)}$ There is a cobra in $H_+$ of definite weight based upon  a
ray $\RR_i(0)$;

\ssk \nin $\mathrm {(ii)}  $  $T|\RR(1) =\id$ (so all the leaves of the ``limiting lamination'' are closed). 
\end{lem}

\begin{proof}
  If there is a cobra of definite weight based upon $\RR(1)$ then
  one of its lifts is the desired cobra based upon $\RR(0)$.  

Otherwise $\eps(x)=1$ a.e. for the return map $T_1$  to $\RR(1)$, so it is
monotonic. Since it is measure preserving, it is identical. 
\end{proof} 

\msk
\subsubsection{Quadratic-like restrictions}
Below, we will restric ourselves to the case of two eyes
(of period $2p+1$). 

Let us consider a cobra segment $S$ based on sone $\RR(n)$ with $n<p$
that begins by crossing  the edge $e_{n-1} $  connecting $K_n$ to $K_{n-1}$.
Assume it is part of some canonical arc $\alpha$.  Then typically it cannot cross the
path $\de_{n-2}$ of edges $e_{j-1}$ consecutively connecting $K_j$ to $K_{j-1}$,  $j=n-2,\dots
1$, for it would lead to an interection between $\alpha$ and $f^* \alpha$. Hence $S$
goes around  $\de_{n-2}$  returning to $\RR(n)$ on the other side of $\de_{n-2}$, or else
it comes back through $e_n$.  
%Let us consider a segment of the first kind. 

\note{How about returning through $e_{n-1}$? }


There are only bounded number of segments of this kind. \note{why?}  
Assume the total weight of the corresponding arcs is of order
$p\WW_\loc$. Then we can push-forward such a segment to obtain at
most $O(p)$ segments of weight of order $\WW_\loc$. These segments
represent cobras based on $\RR_{n+1}$ of the same type, and there are
only a bounded number  of those,   \note{why?} 
so in fact, there is only bounded
number of push-forwards in question.  \note{elaborate!}

Doing this finitely many times,  we obtain a segment $S$  whose
push-forward $f_*S$ encloses $S$, creating a quadratic-like
restriction of the map.  \note{why?} 

\msk
Somewhat different approach. 
Let us  have a cobra $S$ (which is part of a canonical arc $\alpha$) 
 based upon some ray $\Ray(n)$ that together
bound a disk $D$. 
Assume:

\ssk \nin (0)
  $D$ does not contain $K_1 $; 

\ssk \nin (i)
 $\alpha$ is pushable forward 
(as we know, a typical arc is such);  

\ssk\nin (ii)
 $\Ray(n+1)$ is disjoint from $D$;

\ssk\nin (iii) 
There is  a maximal string of at least two consecutive  little Julia sets 
$K_l , \dots , K_m$ contained in $D$, with $m< n$. 

\ssk
Then $f_* (D)$ 
is a cobra-disk  based upon $\Ray(n+1)$ 
that encloses $K_{n+1}$. Hence it cannot be disjoint from $D$.
By (ii), $f_*(D)$ cannot be contained in $D$. 
Hence 

\ssk\nin
a) either $f_*(S)$ enters $D$ through the mouth of $S$ creating a new
cobra on $\Ray(n)$, or else  

\ssk \nin
b)  $f_*(D)\supset D$. 

\ssk\nin 
Assume in addition that

\ssk\nin (iv) 
 $D\supset K_0$

\ssk
 Then Case a) is excuded. Indeed $ f(D) \supset K_1$, 
while $f(D) \sm D \not\supset K_1$. \note{why?} 
Hence $D\supset K_1$ contradicting assumption (0).  

Case  b)  is excluded since otherwise we would obtain 
a  %  (homotopic) 
quadratic-like map $f: D\ra f  (D) $.

\ssk 
How to excule the situaton when $D\not\supset K_0$ (to get rid of
(iv))? In this case, $D$ is sometimes a Thirston obstruction.
Namley, in case a) if $f(D)\sm D $ is 

 

\msk
\subsubsection{Cobras with long tails} 

Let us consider a cobra based on $\RR(1)$ whose tail extends further
into the translation region crossing consecutively $\RR(1), \dots , \RR(N)$ so that the tail region
 bounded by $\RR(1)$ and $\RR(N)$ is empty (no little Julia sets inside). 
Let us pull it back starting with   the shift at the translation region. These
pullbacks contain a bounded number of little Julia sets located outside
the translation region.  Hence, after a bounded number of pullbacks,
we will either lose all these Julia sets, or will stabilize. In the
former case, we obtain an empty cobra that crosses about $N$ rays.
So, if the total weight of such cobras is of order $p\WW_\loc$, then the
total flux they produce is of order $Np \WW_\loc$.  This is not
allowed if $N$ is sufficiently big. 

In the latter case, we obtain a Thurston obstruction.   
 
\sss{Confined situation}
Assume we have cobras encircling  almost the whole set
of little Julia sets $J_{k+p}$, $0\leq k\leq p$. 
Then canonical arcs that start at $J_l$, $0\leq l\leq p$, 
either produce a quadrtic-like restriction, or empty cobras as above, 
or else they are confined. 


\subsection{No closed curves}

\begin{lem}
  If $f$ is not intermediately renormalizable then the $2$-Canonical Closed Curve
  Diagram $\CC$ is empty. 
\end{lem}

\begin{proof}
First note that a curve $\gamma\in \CC$ cannot suround a single Julia set $J_s$ since in this case
the local weight of $J_s$ would be less than $1$. 


 Let us take a curve $\gamma\in \CC$ surrounding the
  smallest number of little Julia sets.  
  Let $s\in [0,p)$ and $t\in (0,p)$ be the smallest ones such that $K_s$ and $K_{-t}$
  are surrounded by $\gamma_s:= \gamma$.
Let us pull $\gamma_s$ back $s+t$ times along the
  orbit of $K_s$. We obtain curves $\gamma_{s-1}, \dots, \gamma_{-t}$. 
By Lemma \ref{flux vs weight},
$$
   \WW(\gamma_{s-n}) \geq d_n^{-1}  \WW(\gamma_{s-n+1}),
$$
where $d_s=2$ and  $d_n=1 $ for $n\not=s$. It follows that all the curves $\gamma_{s-n}$,
belong to $\CC$.
Hence they are disjoint, so the correspoding Jordan disks $D_{s-n}$ are either disjoint or nested.  

Furthermore, since $\gamma$ surrounds the minimal number of little Julia sets,
all these curves 
surround the same number of little Julia sets.
It follows that if the disks $D_{s-n}$ and $D_{s-m}$ are nested then 
$\gamma_{s-n}$ and $\gamma_{s-m}$ are homotopic in $\C\sm \KK$.

Putting the above ingredients together, we see that $D_{-t}\cap \KK$
overlaps with $D\cap \KK$, the map $f^{s+t} : D_{-t}\ra D$ is a branched covering of degree two (if $s>0$) or one
(if $s=0$). 
Moreover, $\gamma_{-t}$ is homotopic to $\gamma$ in $\C\sm \KK$. 
In the former case,  $f$ is intermediately renormalizable;
in the latter case these curves form a Thurston obstruction (that does not exist in our situation).   
\end{proof}

\subsection{Long rectangles} 

\bignote{Written Feb 2017} 

We assume that the period is $2p+1$ (so it is really the doube eye case).

  Let us consider a non-confined arc $\alpha$ that begins and ends deep inside the
  translation region.  Then we have an initial and the final pieces of
  $\alpha$, $\alpha_i$ and $\alpha_f$. Both of them must exit the
  translation region.  Assume they  cross $\Ray (1)$ before exiting the
  region. (It should be possible to achieve  tihs  typically by taking
  $f_*(\alpha)$.) If there is another ray, $\Ray(n)$,  deep inside the tranlation
  region, then pieces of $\Ray(1)$ and $\Ray(n)$, together with pieces
  $A_i$ and $A_f$
  of  $\alpha_i $ and $\alpha_f$, form a long rectangle $R$. 
Let us analyze the structure of this rectangle. 

Assume first that both pieces $A_i $ and $A_f$  cross $\Ray(1)$
positively (``from inside'').  Let us also assume that we have just
one loop on the other side of $\Ray(1)$ and this loop does not contain
$K_0$, so this loop is a  is a cobra. 
Then consider several cases:

1) The rectangle $R$  is empty, i.e., it does not contain any little $K$-sets.
Then our cobra has a long tail.

2)  $R$ contains a string of consequtive little $K$-sets in the
adjunct orbit (of $K_{p+1}$). and then ``dive down " or ``up".
 Let us push forward $\alpha$. Then both horizontal sides of $R$
 go ``up" or ``down'', which forces $f(R)$ and $R$ cross. 

3) The pieces $\alpha_i$ and $\alpha_f$ don't dive, 
  but eventuallay land at the adjunct Julia set.
  Then look at the main string of little Julia sets (the translation orbit of $K_1$)
  and argue that they are confined.

  ***************  } 